% Method -- paper plan Task 4.
% Cross-checked against prototype/healgen/{claims,risk,sufficiency,graph,
% repair,systems,judge}.py at writing time. Describes the code as built:
%  - draft-first decomposition (ledger R2)
%  - 6 CRE features in FEATURE_NAMES order (risk.py)
%  - repair actions actually implemented: rephrase / reduce_confidence /
%    express_uncertainty (+ supported no-op); retrieval is conditional at
%    the system level, not a repair action.
%  - hybrid support rule (ledger R8) documented explicitly.

\section{Method}
\label{sec:method}

\subsection{Problem setting}
Given a question $q$ and a generator $G$, produce a response $r$ together
with a per-claim reliability vector. The system may consult a corpus $C$
through a BM25 retriever, at a cost we treat as a first-class quantity. No
training of $G$ is performed; everything downstream of $G$ is training-free
except the small logistic risk model, which is fit once on a split disjoint
from all evaluation data.

\subsection{Draft and claim decomposition}
$G$ first answers $q$ with no context, yielding a draft. The draft is split
into atomic claims by rule-based segmentation: sentence splitting with
abbreviation guards, conjunctive-clause splitting at \emph{and/but/while/
whereas}, and pronoun resolution that copies the subject onto clauses that
begin with a pronoun, so each claim is a standalone declarative statement
(\texttt{claims.py}). Decomposition is \emph{draft-first}: it operates on
what the model chose to assert, not on the question, so that risk and
sufficiency are properties of asserted content.

\subsection{Claim Risk Estimator (CRE)}
Each claim $c$ is mapped to a six-feature vector
(\texttt{risk.py}): claim length in words; count of non-initial tokens
that contain a digit or begin with a capital letter; maximum BM25 score
of the claim against the corpus; MiniLM cosine
between claim and question; \emph{novelty}, defined as one minus the maximum
MiniLM cosine between the claim and the top-$k$ retrieved passages; and a
rarity flag for all-caps tokens. A logistic regression over standardized
features maps the vector to a risk score in $(0,1)$; the model is stored as
plain JSON coefficients (\texttt{cre\_model.json}), not a pickle. Training
labels come from judge-vs-gold judgments on claims sampled from splits
disjoint from dev/test/judge sets. We report the estimator's training
accuracy against the majority baseline and discuss its degeneracy openly in
\S\ref{sec:limitations}: on this benchmark most draft claims are initially
unsupported, so the positive rate is high and the estimator's practical
value is in \emph{ranking} claims, not in separating classes.

\subsection{Adaptive Evidence Sufficiency Predictor (AESP)}
Sufficiency is embedding-based: $\mathrm{score}(c, E) = \max_{e \in E}
\cos(\mathrm{MiniLM}(c), \mathrm{MiniLM}(e))$, and the claim is sufficient
when the score meets or exceeds a threshold $\tau$ (\texttt{sufficiency.py}). $\tau$
is tuned over $\{0.3, 0.5, 0.7\}$ on the dev split by F1 against judge
labels; the sensitivity of all three settings is reported
(Table~\ref{tab:tau}).  The committed benchmark run used $\tau{=}0.7$,
the value selected by an earlier partial sweep; re-running the full
\NTauPairs-pair sweep (Table~\ref{tab:tau}) selects $0.3$.  On this benchmark
the difference is provably inert: every draft decomposes to a single
claim, which is tested for sufficiency against an initially \emph{empty}
evidence pool, so the pre-retrieval sufficiency score is $0.0 < \tau$
regardless of $\tau$ and retrieval is always issued, while the
post-retrieval sufficiency flag is recorded but never gates any later
decision --- consistent with the no-AESP ablation reproducing the full
system exactly (\S\ref{sec:results}). This judges
whether available evidence is \emph{enough}, which is a different question
from whether it is \emph{relevant} \cite{qiu2026surerag}.

\subsection{Conditional retrieval}
The claim loop is: compute risk; compute sufficiency against the evidence
gathered so far; if and only if the claim is insufficient, issue one BM25
query (Okapi, $k_1{=}1.5$, $b{=}0.75$, top-$k{=}3$) and merge the passages
into the shared pool. Sufficient claims trigger no lookup. The number of
BM25 calls is recorded per response as the retrieval-cost metric, which
makes the loop's central efficiency hypothesis directly testable against a
fixed-retrieval RAG baseline that always queries.

\subsection{Dynamic Claim--Evidence Graph (DCEG)}
Claims and passages become nodes of a bipartite graph (\texttt{graph.py}).
Each claim--passage pair is scored by the judge and stored as a weighted
edge; a claim is \emph{unsupported} when it has no edge above the support
threshold or when it participates in a contradiction edge. The graph
exports to DOT/PNG for inspection; Figure~\ref{fig:dceg} is produced by the
prototype itself.

\subsection{Localized Self-Healing Engine (LSHE)}
Repair operates at sentence granularity. For every unsupported claim, the
sentence that produced it is located by offsets; supported claims are never
touched. The action sequence per unsupported sentence is:
\begin{enumerate}
  \item \textbf{rephrase} --- regenerate \emph{only that sentence} with the
        evidence pool in the prompt; if the regeneration differs from the
        original, splice it in;
  \item \textbf{reduce\_confidence} --- if regeneration yields nothing new,
        append an explicit hedge to that sentence;
  \item \textbf{express\_uncertainty} --- when no evidence is available,
        hedge without grounding.
\end{enumerate}
Sentences not selected for repair are reassembled byte-identically; a
regression test pins this (\texttt{test\_repair\_preserves\_untouched\_
sentences}). The invariants to state plainly: retrieval happens at the
system level when a claim is insufficient, and there is no
\emph{reconnect/drop} action in this prototype --- flagged claims that
survive repair are emitted hedged, not silently removed.

\subsection{Judge and evidence pools}
Support decisions come from an NLI model
(\texttt{DeBERTa-v3-base-mnli-fever-anli}) under a \textbf{hybrid support
rule}: a claim counts as supported if its normalized form (lowercased,
punctuation-stripped, whitespace-collapsed) is a substring of a normalized
passage --- scored $1.0$ --- or if the NLI argmax label on the best passage
is \emph{entailment}; the reported score is the maximum of the two
components, with the NLI softmax taken at unit temperature (no temperature
scaling is fitted anywhere; the judge gate below is also measured at
$T{=}1$).
Evaluation never scores a system against its own retrievals: each question
gets one \emph{frozen shared pool} (gold knowledge plus the union of all
four systems' passages), so no system can win by retrieving more. The judge
itself is validated on labeled contrast pairs before any benchmark run and
its accuracy/AUC are reported as a gate (\S\ref{sec:experiments}). The
lexical branch is a deliberate robustness choice for a small NLI model, and
its risk --- over-crediting lexical overlap --- is acknowledged in
\S\ref{sec:limitations} with a pure-NLI sensitivity discussion.

\subsection{Baselines}
All four systems share the same generator, decoding (greedy,
\texttt{max\_new\_tokens}=64), and BM25 index:
\textbf{Vanilla} answers without retrieval;
\textbf{RAG} always retrieves top-$k$ passages and conditions on them;
\textbf{CoVe} drafts, verifies every claim, and regenerates the whole
response once if anything is unsupported \cite{dhuliawala2023cove};
\textbf{HealGen} runs the full loop above.
